\label{kap:prostor}
\mathquote{Geometry is knowledge of the eternally existent.}{Platón}

V kapitole~\ref{kap:stereometrie} jsme řešili prostorové úlohy tak, že jsme hledali vhodné pravoúhlé trojúhelníky. To funguje, ale vyžaduje to dobrou prostorovou představivost a trochu štěstí. Analytická geometrie v prostoru nabízí jinou cestu: zavedeme souřadnice a všechno spočítáme. Žádné hledání trojúhelníků, žádné kreslení. Jen vektory a rovnice. Pro ty z nás, kdo si v hlavě neumí otočit ani židli, je to vysvobození.

\section{Souřadnice a vektory v prostoru}

\begin{definitionbox}
  \textbf{Kartézská soustava souřadnic} v prostoru má tři navzájem kolmé osy $x$, $y$, $z$. Bod je určen trojicí $A = [a_1, a_2, a_3]$, vektor trojicí $\vec{u} = (u_1, u_2, u_3)$. Všechny vzorce z roviny se přenesou přidáním třetí souřadnice:
  \begin{align*}
    |AB| &= \sqrt{{(b_1 - a_1)}^2 + {(b_2 - a_2)}^2 + {(b_3 - a_3)}^2}, & S_{AB} &= \left[\tfrac{a_1 + b_1}{2}, \tfrac{a_2 + b_2}{2}, \tfrac{a_3 + b_3}{2}\right],\\
    |\vec{u}| &= \sqrt{u_1^2 + u_2^2 + u_3^2}, & \vec{u}\cdot\vec{v} &= u_1v_1 + u_2v_2 + u_3v_3.
  \end{align*}
  I zde platí $\vec{u}\cdot\vec{v} = |\vec{u}||\vec{v}|\cos\varphi$ a vektory jsou kolmé, právě když $\vec{u}\cdot\vec{v} = 0$.
\end{definitionbox}

Jak si bod v prostoru představit? Z počátku ujdeme $a_1$ ve směru osy $x$, pak $a_2$ ve směru osy $y$ a nakonec vystoupáme o $a_3$ nahoru. Bod je roh `krabice' o rozměrech $a_1\times a_2\times a_3$ a délka vektoru $\overrightarrow{OA}$ je její tělesová úhlopříčka. Proto vzorec pro délku vypadá jako Pythagorova věta použitá dvakrát: nejdřív úhlopříčka podstavy $\sqrt{a_1^2 + a_2^2}$, pak s výškou $a_3$.

\begin{center}
  \begin{tikzpicture}[krychle, scale=1.1]
    \draw[->] (0,0,0) -- (3.3,0,0) node[right] {\small $x$};
    \draw[->] (0,0,0) -- (0,4.2,0) node[above right] {\small $y$};
    \draw[->] (0,0,0) -- (0,0,2.9) node[above] {\small $z$};
    \coordinate (A) at (2,3,2);
    \draw[dashed, gray] (2,0,0) -- (2,3,0) -- (0,3,0);
    \draw[dashed, gray] (2,0,0) -- (2,0,2) (2,3,0) -- (2,3,2) (0,3,0) -- (0,3,2);
    \draw[dashed, gray] (0,0,2) -- (2,0,2) -- (2,3,2) -- (0,3,2) -- cycle;
    \draw[dotted, thick] (0,0,0) -- (2,3,0) node[midway, below right, font=\scriptsize] {$\sqrt{a_1^2 + a_2^2}$};
    \draw[->, thick, colorofwisdom] (0,0,0) -- (A);
    \fill (A) circle (1.5pt) node[above right] {\small $A = [2, 3, 2]$};
    \node[below left] at (0,0,0) {\small $O$};
    \node[below] at (2,0,0) {\scriptsize $a_1 = 2$};
    \node[right] at (0,3,0) {\scriptsize $a_2 = 3$};
    \node[left] at (0,0,2) {\scriptsize $a_3 = 2$};
  \end{tikzpicture}
\end{center}

\begin{example}
  Pro $A = [1, 2, 3]$ a $B = [3, -1, 9]$ je $\overrightarrow{AB} = (2, -3, 6)$, $|AB| = \sqrt{4 + 9 + 36} = 7$ a $S_{AB} = \left[2, \frac{1}{2}, 6\right]$.
\end{example}

\subsection{Vektorový součin}

V prostoru existuje ještě jeden součin vektorů, který v rovině nemá obdobu. Výsledkem není číslo, ale vektor.

\begin{definitionbox}
  \textbf{Vektorový součin} vektorů $\vec{u} = (u_1, u_2, u_3)$ a $\vec{v} = (v_1, v_2, v_3)$ je vektor
  \begin{equation*}
    \vec{u}\times\vec{v} = (u_2v_3 - u_3v_2,\ u_3v_1 - u_1v_3,\ u_1v_2 - u_2v_1).
  \end{equation*}
  \textbf{Vlastnosti}:
  \begin{itemize}
    \item $\vec{u}\times\vec{v}$ je \textbf{kolmý} na $\vec{u}$ i na $\vec{v}$,
    \item $|\vec{u}\times\vec{v}| = |\vec{u}||\vec{v}|\sin\varphi$, což je \textbf{obsah rovnoběžníku} určeného vektory $\vec{u}$, $\vec{v}$,
    \item $\vec{u}\times\vec{v} = -\vec{v}\times\vec{u}$ (není komutativní!),
    \item $\vec{u}\times\vec{v} = \vec{o}$, právě když jsou $\vec{u}$ a $\vec{v}$ kolineární.
  \end{itemize}
\end{definitionbox}


\begin{center}
  \begin{tikzpicture}[krychle, scale=1.25]
    \coordinate (O) at (0,0,0); \coordinate (U) at (2.5,0,0); \coordinate (V) at (0.8,1.8,0); \coordinate (UV) at (3.3,1.8,0);
    \draw[->, thick, gray] (O) -- (0,0,-1.1) node[below, font=\small] {$\vec{v}\times\vec{u}$};
    \fill[colorofwisdom!20] (O) -- (U) -- (UV) -- (V) -- cycle;
    \draw[dashed] (U) -- (UV) -- (V);
    \draw[->, thick] (O) -- (U) node[below right] {\small $\vec{u}$};
    \draw[->, thick] (O) -- (V) node[left] {\small $\vec{v}$};
    \node[font=\scriptsize, align=center] at (1.75,0.9,0) {obsah\\$|\vec{u}\times\vec{v}|$};
    \draw[->, thick, colorofwisdom] (O) -- (0,0,2.1) node[above, font=\small] {$\vec{u}\times\vec{v}$};
    \draw (0.25,0,0) -- (0.25,0,0.25) -- (0,0,0.25);
  \end{tikzpicture}
\end{center}
Na obrázku vidíte všechny vlastnosti najednou: výsledek trčí kolmo z roviny, ve které leží $\vec{u}$ a $\vec{v}$, jeho délka je obsah vybarveného rovnoběžníku a po prohození pořadí míří opačně (šedá šipka). Když jsou $\vec{u}$ a $\vec{v}$ rovnoběžné, rovnoběžník se `splácne' na úsečku s nulovým obsahem.

Vzorec vypadá hrozivě, ale dá se zapamatovat jako `determinant'. Každá složka je determinant $2\times 2$ z ostatních dvou souřadnic: první složka vznikne zakrytím prvního sloupce, druhá zakrytím druhého (s obráceným znaménkem), třetí zakrytím třetího:
\begin{equation*}
  \vec{u}\times\vec{v} = \left(\begin{vmatrix} u_2 & u_3\\ v_2 & v_3 \end{vmatrix},\ -\begin{vmatrix} u_1 & u_3\\ v_1 & v_3 \end{vmatrix},\ \begin{vmatrix} u_1 & u_2\\ v_1 & v_2 \end{vmatrix}\right).
\end{equation*}
Kolmost si ověříte přímo: $\vec{u}\cdot(\vec{u}\times\vec{v}) = u_1(u_2v_3 - u_3v_2) + u_2(u_3v_1 - u_1v_3) + u_3(u_1v_2 - u_2v_1) = 0$, protože se všechny členy po dvou vyruší. Směr výsledku určuje \textbf{pravidlo pravé ruky}: pokud prsty pravé ruky ukazují od $\vec{u}$ k $\vec{v}$, palec ukazuje směr $\vec{u}\times\vec{v}$. Fyzici tohle pravidlo používají u magnetické síly a točivého momentu tak často, že je u zkoušek z fyziky vidět sál plný studentů, kteří kroutí pravou rukou ve vzduchu.

\begin{example}
  Pro $\vec{u} = (1, 2, 3)$ a $\vec{v} = (4, 5, 6)$:
  \begin{equation*}
    \vec{u}\times\vec{v} = (2\cdot 6 - 3\cdot 5,\ 3\cdot 4 - 1\cdot 6,\ 1\cdot 5 - 2\cdot 4) = (-3, 6, -3).
  \end{equation*}
  Kontrola: $(1, 2, 3)\cdot(-3, 6, -3) = -3 + 12 - 9 = 0$ a $(4, 5, 6)\cdot(-3, 6, -3) = -12 + 30 - 18 = 0$.
\end{example}

\section{Přímka a rovina v prostoru}

\begin{definitionbox}
  \textbf{Přímka} procházející bodem $A$ se směrovým vektorem $\vec{u}$ má \textbf{parametrické vyjádření}
  \begin{equation*}
    X = A + t\vec{u},\quad t\in\mathbb{R},\qquad\text{tedy}\qquad x = a_1 + tu_1,\quad y = a_2 + tu_2,\quad z = a_3 + tu_3.
  \end{equation*}
  Přímka v prostoru \textbf{nemá} obecnou rovnici. Jedna lineární rovnice v prostoru totiž popisuje rovinu, ne přímku. Přímku lze popsat nejvýš jako průsečnici dvou rovin.

  \vspace{0.2cm}

  \textbf{Rovina} procházející bodem $A$ a určená dvěma nekolineárními vektory $\vec{u}$, $\vec{v}$ má \textbf{parametrické vyjádření}
  \begin{equation*}
    X = A + t\vec{u} + s\vec{v},\quad t, s\in\mathbb{R}.
  \end{equation*}
  \textbf{Obecná rovnice roviny} s normálovým vektorem $\vec{n} = (a, b, c)$ je
  \begin{equation*}
    ax + by + cz + d = 0.
  \end{equation*}
  Normálový vektor roviny určené vektory $\vec{u}$, $\vec{v}$ najdeme jako $\vec{n} = \vec{u}\times\vec{v}$.
\end{definitionbox}


\begin{center}
  \begin{tikzpicture}[scale=1.1]
    \begin{scope}[krychle]
      \draw[thick] (-0.94,-0.5,-0.4) -- (2.84,1.6,1.65) node[right] {\small $p$};
      \draw[->, very thick, colorofwisdom] (0.5,0.3,0.4) -- (1.4,0.8,0.9) node[midway, above left] {\small $\vec{u}$};
      \foreach \t in {-1,0,1,2} {
        \fill ({0.5+0.9*\t},{0.3+0.5*\t},{0.4+0.5*\t}) circle (1.5pt) node[below right, font=\scriptsize] {$t = \t$};
      }
      \node[above left] at (0.5,0.3,0.4) {\small $A$};
    \end{scope}
    \begin{scope}[xshift=5.2cm, krychle]
      \fill[colorofwisdom!15] (-0.8,-0.8,-0.4) -- (2.4,-0.8,0.24) -- (2.4,2.4,1.2) -- (-0.8,2.4,0.56) -- cycle;
      \draw[gray] (-0.8,-0.8,-0.4) -- (2.4,-0.8,0.24) -- (2.4,2.4,1.2) -- (-0.8,2.4,0.56) -- cycle;
      \draw[->, thick] (0,0,0) -- (2,0,0.4) node[below right] {\small $\vec{u}$};
      \draw[->, thick] (0,0,0) -- (0,2,0.6) node[above left] {\small $\vec{v}$};
      \draw[dashed] (0,0,0) -- (1.4,0,0.28) -- (1.4,1,0.58);
      \fill (1.4,1,0.58) circle (1.5pt) node[above right, font=\scriptsize] {$X = A + t\vec{u} + s\vec{v}$};
      \draw[->, thick, colorofwisdom] (0,0,0) -- (-0.24,-0.36,1.2) node[above] {\small $\vec{n} = \vec{u}\times\vec{v}$};
      \fill (0,0,0) circle (1.5pt) node[below left] {\small $A$};
      \node at (2.2,2.2,1.1) [below left] {\small $\rho$};
    \end{scope}
  \end{tikzpicture}
\end{center}
Parametr $t$ funguje jako `čas': v čase $t = 0$ stojíme v bodě $A$ a za každou jednotku času ujdeme vektor $\vec{u}$. Záporné $t$ znamená cestu pozpátku. U roviny máme dva nezávislé `ovladače' $t$ a $s$, a tedy dva rozměry, ve kterých se můžeme pohybovat. Normálový vektor $\vec{n}$ trčí z roviny kolmo jako stožár z podlahy.

Proč má rovina obecnou rovnici a přímka ne? Rovnice $\vec{n}\cdot\overrightarrow{AX} = 0$ říká, že vektor $\overrightarrow{AX}$ je kolmý na $\vec{n}$. V rovině jsou všechny takové body $X$ na jedné přímce. V prostoru tvoří všechny vektory kolmé na $\vec{n}$ celou rovinu. Jedna podmínka v prostoru `ubere' jen jeden rozměr ze tří.

\begin{example}
  Napišme obecnou rovnici roviny $ABC$, kde $A = [1, 0, 0]$, $B = [0, 2, 0]$ a $C = [0, 0, 3]$.

  Vektory $\vec{u} = \overrightarrow{AB} = (-1, 2, 0)$ a $\vec{v} = \overrightarrow{AC} = (-1, 0, 3)$. Normálový vektor
  \begin{equation*}
    \vec{n} = \vec{u}\times\vec{v} = (2\cdot 3 - 0\cdot 0,\ 0\cdot(-1) - (-1)\cdot 3,\ (-1)\cdot 0 - 2\cdot(-1)) = (6, 3, 2).
  \end{equation*}
  Rovnice $6x + 3y + 2z + d = 0$, dosazením $A$: $d = -6$. Tedy $6x + 3y + 2z - 6 = 0$. Kontrola s $B$: $6 - 6 = 0$, s $C$: $6 - 6 = 0$.

  Po vydělení šesti dostaneme hezký tvar $\frac{x}{1} + \frac{y}{2} + \frac{z}{3} = 1$, tzv. \textbf{úsekový tvar}, ze kterého jsou průsečíky s osami vidět okamžitě. A obsah trojúhelníku $ABC$ je polovina obsahu rovnoběžníku: $S = \frac{1}{2}|\vec{n}| = \frac{1}{2}\sqrt{36 + 9 + 4} = \frac{7}{2}$.
\end{example}

\begin{center}
  \begin{tikzpicture}[krychle, scale=1.15]
    \draw[->] (0,0,0) -- (1.8,0,0) node[right] {\small $x$};
    \draw[->] (0,0,0) -- (0,2.9,0) node[above right] {\small $y$};
    \draw[->] (0,0,0) -- (0,0,3.6) node[above] {\small $z$};
    \fill[colorofwisdom!30, opacity=0.8] (1,0,0) -- (0,2,0) -- (0,0,3) -- cycle;
    \draw[thick, colorofwisdom] (1,0,0) -- (0,2,0) -- (0,0,3) -- cycle;
    \fill (1,0,0) circle (1.5pt) node[below] {\small $A = [1, 0, 0]$};
    \fill (0,2,0) circle (1.5pt) node[right] {\small $B = [0, 2, 0]$};
    \fill (0,0,3) circle (1.5pt) node[left] {\small $C = [0, 0, 3]$};
    \node[font=\small] at (2.6,0,2.2) {$\dfrac{x}{1} + \dfrac{y}{2} + \dfrac{z}{3} = 1$};
  \end{tikzpicture}
\end{center}
Rovina `odřízne' z prvního oktantu trojúhelníkový kousek. Čísla ve jmenovatelích úsekového tvaru jsou přesně místa, kde rovina protíná osy.

\section{Vzájemné polohy}

\begin{definitionbox}
  \textbf{Dvě přímky} $p: X = A + t\vec{u}$ a $q: X = B + s\vec{v}$:
  \begin{itemize}
    \item pokud jsou $\vec{u}$ a $\vec{v}$ kolineární, jsou přímky \textbf{rovnoběžné} (různé, nebo totožné, pokud $B\in p$),
    \item jinak řešíme soustavu $A + t\vec{u} = B + s\vec{v}$ (tři rovnice pro dvě neznámé). Pokud má řešení, jsou \textbf{různoběžné}, pokud nemá, jsou \textbf{mimoběžné}.
  \end{itemize}

  \textbf{Přímka} $p: X = A + t\vec{u}$ a \textbf{rovina} $\rho: \vec{n}\cdot X + d = 0$:
  \begin{itemize}
    \item pokud $\vec{n}\cdot\vec{u}\neq 0$, jsou \textbf{různoběžné}. Průsečík najdeme dosazením parametrického vyjádření přímky do rovnice roviny,
    \item pokud $\vec{n}\cdot\vec{u} = 0$, je přímka s rovinou \textbf{rovnoběžná}, nebo v ní \textbf{leží} (pokud $A\in\rho$).
  \end{itemize}

  \textbf{Dvě roviny}: jsou \textbf{rovnoběžné}, pokud mají kolineární normálové vektory (a totožné, pokud jsou rovnice násobkem jedna druhé). Jinak jsou \textbf{různoběžné} a průsečnici najdeme vyřešením soustavy dvou rovnic se třemi neznámými (jedna neznámá bude parametr).
\end{definitionbox}

\begin{center}
  \begin{tikzpicture}[scale=1]
    \begin{scope}[krychle]
      \draw[thick] (0,0,0) -- (2,1,1) node[right] {\small $p$};
      \draw[thick] (2,0,0) -- (0,1,1) node[left] {\small $q$};
      \fill (1,0.5,0.5) circle (1.6pt);
      \node[font=\small] at (1,0,-0.6) {různoběžné};
    \end{scope}
    \begin{scope}[xshift=4cm, krychle]
      \draw[thick] (0,0,0) -- (2,0.5,1) node[right] {\small $p$};
      \draw[thick] (0,1,0.3) -- (2,1.5,1.3) node[right] {\small $q$};
      \node[font=\small] at (1,0,-0.6) {rovnoběžné};
    \end{scope}
    \begin{scope}[xshift=8cm, krychle]
      \draw[gray!60, dashed] (0,0,0) -- (2,0,0) -- (2,1.5,0) -- (0,1.5,0) -- cycle;
      \draw[gray!60, dashed] (0,0,1.2) -- (2,0,1.2) -- (2,1.5,1.2) -- (0,1.5,1.2) -- cycle;
      \draw[gray!60, dashed] (0,0,0) -- (0,0,1.2) (2,0,0) -- (2,0,1.2) (2,1.5,0) -- (2,1.5,1.2) (0,1.5,0) -- (0,1.5,1.2);
      \draw[thick] (-0.3,0.4,0) -- (2.3,0.4,0) node[right] {\small $p$};
      \draw[thick] (1,-0.3,1.2) -- (1,1.9,1.2) node[above right] {\small $q$};
      \node[font=\small] at (1,0,-0.6) {mimoběžné};
    \end{scope}
  \end{tikzpicture}
\end{center}
Mimoběžky jsou to, co v rovině nemáme: dvě přímky, které se nikdy neprotnou, a přesto nejsou rovnoběžné. Představte si přímku vedenou po podlaze podél jedné stěny a přímku po stropě podél stěny sousední. Krabice na obrázku pomáhá: kdyby tam nebyla, na papíře by se přímky zdánlivě křížily. U obrázků prostoru se proto vždy vyplatí přikreslit nějaký `rám'.

\begin{center}
  \begin{tikzpicture}[scale=1]
    \begin{scope}[krychle]
      \fill[colorofwisdom!15] (0,0,0) -- (2.4,0,0) -- (2.4,2,0) -- (0,2,0) -- cycle; \draw[gray] (0,0,0) -- (2.4,0,0) -- (2.4,2,0) -- (0,2,0) -- cycle;
      \draw[thick, dashed] (1.1,0.9,-0.8) -- (1.34,0.98,0);
      \draw[thick] (1.34,0.98,0) -- (1.7,1.1,1.2) node[right] {\small $p$};
      \fill (1.34,0.98,0) circle (1.6pt);
      \node[font=\small] at (1.2,0,-1.1) {různoběžné (průsečík)};
    \end{scope}
    \begin{scope}[xshift=4.3cm, krychle]
      \fill[colorofwisdom!15] (0,0,0) -- (2.4,0,0) -- (2.4,2,0) -- (0,2,0) -- cycle; \draw[gray] (0,0,0) -- (2.4,0,0) -- (2.4,2,0) -- (0,2,0) -- cycle;
      \draw[thick] (0.2,0.3,0.8) -- (2.4,1.5,0.8) node[right] {\small $p$};
      \draw[dashed, gray] (0.2,0.3,0) -- (2.4,1.5,0);
      \draw[dotted] (0.2,0.3,0) -- (0.2,0.3,0.8) (2.4,1.5,0) -- (2.4,1.5,0.8);
      \node[font=\small] at (1.2,0,-1.1) {rovnoběžné};
    \end{scope}
    \begin{scope}[xshift=8.6cm, krychle]
      \fill[colorofwisdom!15] (0,0,0) -- (2.4,0,0) -- (2.4,2,0) -- (0,2,0) -- cycle; \draw[gray] (0,0,0) -- (2.4,0,0) -- (2.4,2,0) -- (0,2,0) -- cycle;
      \draw[thick] (0.2,0.3,0) -- (2.2,1.7,0) node[above] {\small $p$};
      \node[font=\small] at (1.2,0,-1.1) {přímka leží v rovině};
    \end{scope}
  \end{tikzpicture}
\end{center}
Podmínka $\vec{n}\cdot\vec{u} = 0$ znamená, že směr přímky je kolmý na normálu, tedy `vodorovný' vzhledem k rovině. Taková přímka buď letí nad rovinou jako letadlo v konstantní výšce, nebo po ní jede jako auto po silnici.

\begin{example}
  Rozhodněme o poloze přímek $p: X = [1, 0, 0] + t(1, 1, 0)$ a $q: X = [0, 1, 1] + s(1, -1, 0)$. Směrové vektory kolineární nejsou. Soustava:
  \begin{align*}
    1 + t &= s,\\
    t &= 1 - s,\\
    0 &= 1.
  \end{align*}
  Třetí rovnice je spor, takže přímky jsou \textbf{mimoběžné}. Kdyby $q$ procházela bodem $[0, 1, 0]$, třetí rovnice by zněla $0 = 0$ a z prvních dvou bychom dostali $t = 0$, $s = 1$ a průsečík $[1, 0, 0]$.
\end{example}

\begin{example}
  Průsečík přímky $X = [1, 1, 1] + t(1, 2, 3)$ s rovinou $x + y + z = 12$: dosadíme $(1 + t) + (1 + 2t) + (1 + 3t) = 12$, tedy $6t = 9$ a $t = \frac{3}{2}$. Průsečík je $\left[\frac{5}{2}, 4, \frac{11}{2}\right]$.
\end{example}

\begin{example}
  Průsečnice rovin $x + y + z = 3$ a $x - y = 1$: zvolíme $y = t$, pak $x = 1 + t$ a $z = 3 - x - y = 2 - 2t$. Průsečnice je $X = [1, 0, 2] + t(1, 1, -2)$. Kontrola: směrový vektor $(1, 1, -2)$ je kolmý na oba normálové vektory $(1, 1, 1)$ a $(1, -1, 0)$, a tedy je rovnoběžný s $\vec{n}_1\times\vec{n}_2$.
\end{example}

\section{Metrické úlohy}

\begin{definitionbox}
  \textbf{Vzdálenost bodu} $M = [m_1, m_2, m_3]$ od roviny $ax + by + cz + d = 0$:
  \begin{equation*}
    v = \frac{|am_1 + bm_2 + cm_3 + d|}{\sqrt{a^2 + b^2 + c^2}}.
  \end{equation*}
  \textbf{Vzdálenost mimoběžek} $p: X = A + t\vec{u}$ a $q: X = B + s\vec{v}$:
  \begin{equation*}
    v = \frac{|\overrightarrow{AB}\cdot(\vec{u}\times\vec{v})|}{|\vec{u}\times\vec{v}|}.
  \end{equation*}
  \textbf{Odchylky}:
  \begin{align*}
    \text{dvou přímek:}\quad\cos\varphi &= \frac{|\vec{u}\cdot\vec{v}|}{|\vec{u}||\vec{v}|}, &
    \text{dvou rovin:}\quad\cos\varphi &= \frac{|\vec{n}_1\cdot\vec{n}_2|}{|\vec{n}_1||\vec{n}_2|},\\
    \text{přímky a roviny:}\quad\sin\varphi &= \frac{|\vec{u}\cdot\vec{n}|}{|\vec{u}||\vec{n}|}.
  \end{align*}
\end{definitionbox}

Všechny tyto vzorce vypadají jako něco, co je potřeba se naučit nazpaměť. Ve skutečnosti za nimi stojí jediná myšlenka: \textbf{průmět vektoru do vhodného směru}. Pojďme si je projít s obrázky.

\subsection*{Vzdálenost bodu od roviny}

Představte si, že $M$ je lampa visící nad podlahou $\rho$. Nejkratší cesta z lampy na podlahu vede svisle dolů, tedy ve směru normály $\vec{n}$. Vezmeme libovolný bod $P$ podlahy. Vektor $\overrightarrow{PM}$ míří šikmo, ale nás zajímá jen jeho `svislá' složka, tedy jeho průmět do směru $\vec{n}$.

\begin{center}
  \begin{tikzpicture}[krychle, scale=1.15]
    \fill[colorofwisdom!15] (0,0,0) -- (4,0,0) -- (4,3,0) -- (0,3,0) -- cycle;
    \draw[gray] (0,0,0) -- (4,0,0) -- (4,3,0) -- (0,3,0) -- cycle;
    \node at (3.7,2.7,0) {\small $\rho$};
    \coordinate (M) at (2.6,1.6,2); \coordinate (M') at (2.6,1.6,0); \coordinate (P) at (0.8,0.5,0);
    \draw[dashed] (P) -- (M');
    \draw[->, thick] (P) -- (M) node[midway, above left] {\small $\overrightarrow{PM}$};
    \draw[very thick, colorofwisdom] (M) -- (M') node[midway, right] {\small $v$};
    \draw (2.39,1.47,0) -- (2.39,1.47,0.2) -- (2.6,1.6,0.2);
    \draw[->, thick, colorofwisdom] (P) -- (0.8,0.5,1.1) node[left] {\small $\vec{n}$};
    \draw[dotted] (0.8,0.5,1.1) -- (0.8,0.5,2) -- (M);
    \fill (M) circle (1.6pt) node[above right] {\small $M$};
    \fill (M') circle (1.6pt) node[below right] {\small $M'$};
    \fill (P) circle (1.6pt) node[below left] {\small $P$};
  \end{tikzpicture}
\end{center}

Délka průmětu je $|\overrightarrow{PM}|\cdot|\cos\alpha|$, kde $\alpha$ je úhel mezi $\overrightarrow{PM}$ a $\vec{n}$. Ze skalárního součinu $\overrightarrow{PM}\cdot\vec{n} = |\overrightarrow{PM}||\vec{n}|\cos\alpha$ dostaneme
\begin{equation*}
  v = \frac{|\overrightarrow{PM}\cdot\vec{n}|}{|\vec{n}|}.
\end{equation*}
A teď rozepišme čitatel. Pro $P = [p_1, p_2, p_3]$ je
\begin{equation*}
  \overrightarrow{PM}\cdot\vec{n} = a(m_1 - p_1) + b(m_2 - p_2) + c(m_3 - p_3) = am_1 + bm_2 + cm_3 - \underbrace{(ap_1 + bp_2 + cp_3)}_{=\,-d} = am_1 + bm_2 + cm_3 + d,
\end{equation*}
protože $P$ leží v rovině, a tedy $ap_1 + bp_2 + cp_3 + d = 0$. To je celý vzorec. Všimněte si, že nezáleží na tom, který bod $P$ vybereme. Posouváním $P$ po podlaze se mění jen `vodorovná' část vektoru $\overrightarrow{PM}$, svislá zůstává stejná. Mimochodem, výraz $am_1 + bm_2 + cm_3 + d$ bez absolutní hodnoty má ještě jednu informaci: jeho znaménko říká, na \textit{které straně} roviny bod leží.

\subsection*{Vzdálenost mimoběžek}

Každé dvě mimoběžky leží ve dvou rovnoběžných rovinách, jako přímka na podlaze a přímka na stropě. Jejich vzdálenost je tedy prostě `výška místnosti'. Obě roviny mají stejnou normálu, která musí být kolmá na $\vec{u}$ i na $\vec{v}$. Takový vektor už známe: $\vec{u}\times\vec{v}$. Vzdálenost je tak délka průmětu libovolného vektoru $\overrightarrow{AB}$ (z podlahy na strop) do směru $\vec{u}\times\vec{v}$, přesně jako výše.

\begin{center}
  \begin{tikzpicture}[krychle, scale=1.15]
    \fill[colorofwisdom!15] (-0.5,-0.3,0) -- (3.5,-0.3,0) -- (3.5,2.5,0) -- (-0.5,2.5,0) -- cycle;
    \draw[gray] (-0.5,-0.3,0) -- (3.5,-0.3,0) -- (3.5,2.5,0) -- (-0.5,2.5,0) -- cycle;
    \draw[thick] (-0.5,1,0) -- (3.5,1,0) node[right] {\small $p$ (směr $\vec{u}$)};
    \draw[very thick, colorofwisdom] (1.5,1,0) -- (1.5,1,1.6) node[midway, right] {\small $v$};
    \draw[->, thick, colorofwisdom] (0.3,1,0) -- (0.3,1,1) node[left] {\small $\vec{u}\times\vec{v}$};
    \fill[colorofwisdom!8, opacity=0.7] (-0.5,-0.3,1.6) -- (3.5,-0.3,1.6) -- (3.5,2.5,1.6) -- (-0.5,2.5,1.6) -- cycle;
    \draw[gray] (-0.5,-0.3,1.6) -- (3.5,-0.3,1.6) -- (3.5,2.5,1.6) -- (-0.5,2.5,1.6) -- cycle;
    \draw[thick] (1.5,-0.3,1.6) -- (1.5,2.5,1.6) node[above right] {\small $q$ (směr $\vec{v}$)};
    \draw[->, thick] (0.3,1,0) -- (1.5,2.1,1.6);
    \node[right] at (1.1,1.65,1.05) {\small $\overrightarrow{AB}$};
    \fill (0.3,1,0) circle (1.6pt) node[below] {\small $A$};
    \fill (1.5,2.1,1.6) circle (1.6pt) node[above left] {\small $B$};
  \end{tikzpicture}
\end{center}

Vzorec má i krásný geometrický význam. Číslo $\overrightarrow{AB}\cdot(\vec{u}\times\vec{v})$ se jmenuje \textbf{smíšený součin} a jeho absolutní hodnota je \textbf{objem rovnoběžnostěnu} s hranami $\overrightarrow{AB}$, $\vec{u}$, $\vec{v}$. Proč? Objem je podstava krát výška, podstava je rovnoběžník s obsahem $|\vec{u}\times\vec{v}|$ a výška je průmět $\overrightarrow{AB}$ do kolmého směru. Vzorec pro vzdálenost mimoběžek tedy říká
\begin{equation*}
  v = \frac{\text{objem}}{\text{obsah podstavy}} = \text{výška}.
\end{equation*}

\subsection*{Odchylky}

Odchylka dvou rovin se měří jako úhel otevřené knihy: v rovině kolmé na hřbet. Normála každé stránky je od stránky otočená o $90^\circ$, takže když se stránky rozevřou o úhel $\varphi$, normály se rozevřou o stejný úhel. Proto stačí změřit úhel normál. Absolutní hodnota ve vzorci zajistí, že dostaneme úhel mezi $0^\circ$ a $90^\circ$. Normálu totiž můžeme vzít i opačně a odchylka nemá záviset na tom, jak se rozhodneme.

\begin{center}
  \begin{tikzpicture}[scale=1.3]
    \begin{scope}[krychle]
      \fill[colorofwisdom!15] (0,0,0) -- (3,0,0) -- (3,2,0) -- (0,2,0) -- cycle;
      \draw[gray] (0,0,0) -- (3,0,0) -- (3,2,0) -- (0,2,0) -- cycle;
      \coordinate (P) at (1,1,0);
      \coordinate (Q) at ({1+1.9*cos(35)},1,{1.9*sin(35)});
      \coordinate (Q') at ({1+1.9*cos(35)},1,0);
      \coordinate (N) at (1,1,1.4);
      \coordinate (R) at ({1-0.8*cos(35)},1,{-0.8*sin(35)});
    \end{scope}
    \draw[thick, dashed] (R) -- (P);
    \draw[thick] (P) -- (Q) node[right] {\small $p$};
    \draw[dashed] (P) -- (Q');
    \draw[dotted] (Q) -- (Q');
    \draw[->, thick, colorofwisdom] (P) -- (N) node[above] {\small $\vec{n}$};
    \draw ($(P)+(0:0.75)$) arc (0:35:0.75);
    \node at ($(P)+(17:1)$) {\small $\varphi$};
    \draw[colorofwisdom] ($(P)+(35:0.45)$) arc (35:90:0.45);
    \node at ($(P)+(64:0.85)$) {\scriptsize $90^\circ{-}\varphi$};
    \node[font=\small] at (1.8,-0.6) {přímka a rovina};
    \begin{scope}[xshift=5.2cm]
      \begin{scope}[krychle]
        \fill[colorofwisdom!15] (0,0,0) -- (2,0,0) -- (2,2.2,0) -- (0,2.2,0) -- cycle;
        \draw[gray] (0,0,0) -- (2,0,0) -- (2,2.2,0) -- (0,2.2,0) -- cycle;
        \fill[colorofwisdom!30, opacity=0.7] (0,0,0) -- ({2*cos(65)},0,{2*sin(65)}) -- ({2*cos(65)},2.2,{2*sin(65)}) -- (0,2.2,0) -- cycle;
        \draw[gray] (0,0,0) -- ({2*cos(65)},0,{2*sin(65)}) -- ({2*cos(65)},2.2,{2*sin(65)}) -- (0,2.2,0);
        \draw[very thick] (0,-0.2,0) -- (0,2.4,0);
        \coordinate (O) at (0,1.1,0);
        \coordinate (K1) at (1.7,1.1,0);
        \coordinate (K2) at ({1.7*cos(65)},1.1,{1.7*sin(65)});
        \coordinate (N1) at (0,1.1,1.2);
        \coordinate (N2) at ({-1.2*sin(65)},1.1,{1.2*cos(65)});
      \end{scope}
      \draw[thick] (O) -- (K1) (O) -- (K2);
      \draw ($(O)+(0:0.8)$) arc (0:65:0.8);
      \node at ($(O)+(32:1.05)$) {\small $\varphi$};
      \draw[->, thick, colorofwisdom] (O) -- (N1) node[above] {\small $\vec{n}_1$};
      \draw[->, thick, colorofwisdom] (O) -- (N2) node[above left] {\small $\vec{n}_2$};
      \draw[colorofwisdom] ($(O)+(90:0.55)$) arc (90:155:0.55);
      \node[colorofwisdom] at ($(O)+(122:0.8)$) {\small $\varphi$};
      \node[font=\small] at (1.2,-0.6) {dvě roviny};
    \end{scope}
  \end{tikzpicture}
\end{center}

Odchylka dvou přímek je jednodušší: je to úhel jejich směrových vektorů (opět s absolutní hodnotou). U mimoběžek si jednu přímku posuneme tak, aby druhou protínala, a měříme úhel tam.

U odchylky přímky a roviny je sinus, ne kosinus. Proč? Podívejte se na levý obrázek. Úhel přímky a roviny měříme mezi přímkou a jejím `stínem' (kolmým průmětem) v rovině. Vzorec s kosinem by nám ale dal úhel mezi přímkou a \textit{normálou} roviny. Úhel mezi přímkou a rovinou je doplněk do $90^\circ$ a $\cos(90^\circ - \varphi) = \sin\varphi$.

\begin{example}
  Vraťme se ke krychli z kapitoly~\ref{kap:stereometrie} a spočítejme všechno znovu, tentokrát bez hledání trojúhelníků. Krychle s hranou $1$ má vrcholy $A = [0, 0, 0]$, $B = [1, 0, 0]$, $C = [1, 1, 0]$, $D = [0, 1, 0]$, $E = [0, 0, 1]$, $F = [1, 0, 1]$, $G = [1, 1, 1]$, $H = [0, 1, 1]$.
  \begin{itemize}
    \item \textbf{Odchylka tělesové úhlopříčky od podstavy}: $\vec{u} = \overrightarrow{AG} = (1, 1, 1)$, normála podstavy $\vec{n} = (0, 0, 1)$. $\sin\varphi = \frac{1}{\sqrt{3}}$, $\varphi\approx 35{,}26^\circ$. Stejně jako minule.
    \item \textbf{Vzdálenost $A$ od roviny $BDE$}: rovina prochází body $[1, 0, 0]$, $[0, 1, 0]$, $[0, 0, 1]$, má tedy rovnici $x + y + z - 1 = 0$. Vzdálenost $\frac{|-1|}{\sqrt{3}} = \frac{\sqrt{3}}{3}$, tedy přesně třetina tělesové úhlopříčky.
    \item \textbf{Vzdálenost mimoběžek $AB$ a $CG$}: $\vec{u} = (1, 0, 0)$, $\vec{v} = (0, 0, 1)$, $\vec{u}\times\vec{v} = (0, -1, 0)$, $\overrightarrow{AC} = (1, 1, 0)$. Vzdálenost $\frac{|-1|}{1} = 1$.
    \item \textbf{Odchylka rovin $ABG$ a $ABC$}: rovina $ABG$ obsahuje vektory $(1, 0, 0)$ a $(1, 1, 1)$, normála $(0, -1, 1)$. Rovina $ABC$ má normálu $(0, 0, 1)$. $\cos\varphi = \frac{1}{\sqrt{2}}$, $\varphi = 45^\circ$.
  \end{itemize}
  Všechno sedí s výsledky z kapitoly o stereometrii. Kdo nemá prostorovou představivost, má souřadnice.
\end{example}

\begin{center}
  \begin{tikzpicture}[krychle, scale=2.7]
    \krychle
    \fill[colorofwisdom!25, opacity=0.7] (B) -- (D) -- (E) -- cycle;
    \draw[colorofwisdom] (B) -- (E) -- (D);
    \draw[colorofwisdom, dashed] (B) -- (D);
    \coordinate (T) at (1/3,1/3,1/3);
    \draw[very thick, colorofwisdom] (A) -- (T);
    \draw[dashed, colorofwisdom] (T) -- (G);
    \fill (T) circle (0.7pt) node[right, font=\scriptsize] {$T$};
    \node[font=\scriptsize, align=left, anchor=west] at (1.6,0,0.5) {rovina $BDE$ protíná\\úhlopříčku $AG$ v bodě $T$,\\$|AT| = \frac{1}{3}|AG| = \frac{\sqrt{3}}{3}$};
  \end{tikzpicture}
\end{center}
Na obrázku je vidět i to, proč vzdálenost $A$ od roviny $BDE$ měříme právě po úhlopříčce $AG$: ta má směr $(1, 1, 1)$, což je přesně normálový vektor roviny $x + y + z - 1 = 0$. Úhlopříčka tedy rovinu protíná kolmo.

\begin{problem}
  \begin{enumerate}[label=(\alph*)]
    \item Napište obecnou rovnici roviny, která prochází bodem $[2, -1, 3]$ a je kolmá na vektor $(1, 4, -2)$.
    \item Určete vzájemnou polohu přímky $X = [1, 0, 2] + t(2, 1, -1)$ a roviny $x - y + z - 4 = 0$. Pokud jsou rovnoběžné, určete jejich vzdálenost.
    \item Určete vzdálenost rovnoběžných rovin $2x - y + 2z + 3 = 0$ a $2x - y + 2z - 6 = 0$.
    \item Určete odchylku rovin $x + y + z = 1$ a $z = 0$.
    \item Určete obsah trojúhelníku $A = [1, 1, 0]$, $B = [3, 1, 0]$, $C = [1, 4, 4]$.
  \end{enumerate}
\end{problem}

\begin{solution}
  \begin{enumerate}[label=(\alph*)]
    \item $x + 4y - 2z + d = 0$, $2 - 4 - 6 + d = 0$, tedy $x + 4y - 2z + 8 = 0$.
    \item $\vec{n}\cdot\vec{u} = 2 - 1 - 1 = 0$, takže je přímka s rovinou rovnoběžná nebo v ní leží. Bod $[1, 0, 2]$: $1 - 0 + 2 - 4 = -1\neq 0$, rovnoběžná. Vzdálenost $\frac{1}{\sqrt{3}} = \frac{\sqrt{3}}{3}$.
    \item Bod první roviny, například $[0, 3, 0]$, má od druhé vzdálenost $\frac{|-3 - 6|}{\sqrt{4 + 1 + 4}} = \frac{9}{3} = 3$.
    \item $\cos\varphi = \frac{1}{\sqrt{3}}$, $\varphi\approx 54{,}7^\circ$.
    \item $\overrightarrow{AB} = (2, 0, 0)$, $\overrightarrow{AC} = (0, 3, 4)$, $\overrightarrow{AB}\times\overrightarrow{AC} = (0\cdot 4 - 0\cdot 3,\ 0\cdot 0 - 2\cdot 4,\ 2\cdot 3 - 0\cdot 0) = (0, -8, 6)$, $S = \frac{1}{2}\cdot 10 = 5$.
  \end{enumerate}
\end{solution}

\section{Příběh --- Most v Dublinu}

Irský matematik William Rowan Hamilton byl zázračné dítě. Ve třinácti letech prý mluvil třinácti jazyky a ve dvaadvaceti, ještě jako student, byl jmenován královským astronomem Irska. Přes deset let se pak snažil vyřešit problém, který ho pronásledoval: jak `násobit' trojice čísel, tedy body v prostoru, podobně jako komplexní čísla (kapitola~\ref{kap:komplexni}) umožňují násobit body v rovině. Každé ráno se ho prý děti u snídaně ptaly: `Tak co, tati, umíš už násobit trojice?' A on smutně odpovídal, že je umí jen sčítat a odčítat.

Šestnáctého října 1843 šel s manželkou podél kanálu v Dublinu a najednou mu došlo, že trojice nestačí. Potřeboval čtyři čísla. Byl tak nadšený, že nožem vyryl do kamene Broughamova mostu vzorec
\begin{equation*}
  i^2 = j^2 = k^2 = ijk = -1,
\end{equation*}
který definuje takzvané \textbf{kvaterniony}. Nápis už dávno zvětral, ale na mostě je dnes pamětní deska a matematici se k ní každý rok 16. října vydávají na pochod. Kvaterniony se dnes používají k počítání rotací v počítačových hrách, v robotice i v řízení vesmírných sond. Moderní vektorový součin, který jsme poznali v této kapitole, vznikl vlastně jako zjednodušený kus Hamiltonova násobení kvaternionů.

Hamiltonův osobní život tak šťastný nebyl. Celý život miloval ženu jménem Catherine Disneyová, kterou rodina provdala za jiného. V pozdějších letech propadl alkoholu a stal se samotářem. Když po jeho smrti v roce 1865 procházeli jeho pracovnu, našli prý mezi hromadami rukopisů i talíře se zbytky jídla, které tam ležely měsíce, protože odmítal, aby mu kdokoli sahal na papíry.

\section{Shrnutí a otázky k ústní zkoušce}

\begin{problem}
  Odpovězte vlastními slovy:
  \begin{enumerate}
    \item Jak se počítá vzdálenost bodů, velikost vektoru a skalární součin v prostoru?
    \item Co je vektorový součin? Jaké má vlastnosti a k čemu se používá?
    \item Jak vypadá parametrické vyjádření přímky a roviny v prostoru? Proč přímka v prostoru nemá obecnou rovnici?
    \item Jak najdete obecnou rovnici roviny procházející třemi body?
    \item Jak určíte vzájemnou polohu dvou přímek v prostoru? Jak poznáte mimoběžky?
    \item Jak určíte vzájemnou polohu přímky a roviny a dvou rovin?
    \item Jak se počítá vzdálenost bodu od roviny a vzdálenost mimoběžek?
    \item Jak se počítají odchylky přímek a rovin? Proč je u odchylky přímky a roviny sinus?
  \end{enumerate}
\end{problem}
